\section{Scope, Chronology, and Evidence}

\subsection{Identity, names, and range}

Subject: Roberto Francesco Romolo Bellarmino (Latin \textit{Robertus
Bellarminus}; ordinarily Robert Bellarmine in English), 1542--1621:
Jesuit, priest, professor at Louvain and the Roman College,
consultor of the Holy Office, cardinal (1599), archbishop of Capua
(1602--1605), controversial theologian, catechist, spiritual
writer; beatified 1923, canonized 1930, declared Doctor of the
Universal Church 1931. Principal places: Montepulciano, Rome,
Florence, Mondov\`i, Padua, Louvain (with ordination stations at
Li\`ege and Ghent), Paris (legation, 1590), Naples (provincialate),
Capua; Venice and England as fields of controversy only. Whole-life
range secure at day level at both ends: born 4~October 1542, died
17~September 1621, both dates carried by the process-based official
acts and the early record.

\subsection{Reader, question, and thesis boundary}

Written for a serious general reader, student, or researcher who
wants one coherent, evidence-graded life. The governing questions
are stated in the opening section. The organizing synthesis---a
massively documented life whose contested episodes (Sixtine Index,
Bruno, Galileo, the cause's delay) are each governed by a small
number of identifiable documents, around which interpretation has
accreted---is historical. This study does not adjudicate the
theology of grace, the doctrine of papal temporal power, the
Galileo affair's larger science-and-faith questions, or any
doctrinal matter; it claims no theological, canonical,
archival-diplomatic, or rights review and no ecclesiastical
approval. The 1633 trial of Galileo is outside scope except as
reception of the 1616 documents. Aloysius Gonzaga, Bruno, Sarpi,
James~I, and the popes of the period appear only where Bellarmine's
evidence requires them.

\subsection{Corpus and citation conventions}

The autobiography is cited from D\"ollinger--Reusch,
\work{Die Selbstbiographie des Cardinals Bellarmin} (Bonn, 1887),
by page of that edition, read in the public scan; its Latin was
read directly, and renderings are project translations. Galileo
materials are cited from Favaro's \work{Le Opere di Galileo
Galilei, Edizione Nazionale}, vols.~XI--XII and XIX, by document
number and page, read in public scans; Italian and Latin read
directly; project translations marked, except the 5~March 1616
decree, where the public-domain Sturge translation (von Gebler,
1879) was followed and checked. The \work{Disputationes} structure
was verified in the Venice 1599 and Venice 1721 printings (public
scans). Official acts are cited from the \work{Acta Apostolicae
Sedis} by volume and page, read in the Vatican website's scans
(vols.~22, 23, 60). English of James~I and the oath follows
McIlwain's 1918 edition. The 1907--1912 Catholic Encyclopedia
articles are cited by article and year, read at New Advent. All
web witnesses were retrieved and read on 2026-07-25; mutable
facts (calendar, shrine, patronage) carry that cutoff.

\subsection{Evidence classes}

\begin{evidencekey}
\evidenceclass{A}{Subject-authored evidence: the 1613
autobiography (retrospective, edifying, third-person; its genre
limits stated in place); Bellarmine's letters and documents in
Favaro's edition, at the witness level Favaro states (autograph,
signed original, contemporary copy); his published works as read
in named printings.}
\evidenceclass{N}{Contemporary institutional and external
records: the Holy Office and Index documents of 1615--1616 in
Favaro vol.~XIX; the printed 1616 Index decree; the English
statutes, oath, breves, and royal tracts in McIlwain.}
\evidenceclass{E}{Early received tradition: the 1622--1626
process hearings and early lives (Fuligatti 1624 and successors),
used only at reported level through named intermediaries; the
funeral acclamations as reported by the process.}
\evidenceclass{L}{Later devotional and popular tradition:
unverified patronage attributions; the
better-arguments-than-his-opponents quip; post-1923 shrine
devotion.}
\evidenceclass{M}{Official ecclesial reception, each act within
its object: the decrees of 1920--1930; \work{Lux illa} (AAS 22,
593--604); the canonization formula; \work{Providentissimus Deus}
(AAS 23, 433--438); \work{Roma altera} (AAS 60, 12--13); the
current General Roman Calendar; Benedict XVI's catechesis of
23~February 2011.}
\evidenceclass{K}{Named modern critical reconstruction:
D\"ollinger--Reusch (1887) on the autobiography and the cause;
von Gebler (1879) on the 1616 documents; the Catholic
Encyclopedia authors (1907--1912); each dated and characterized.}
\evidenceclass{S}{Source-grounded project synthesis, never
attributed to a source.}
\end{evidencekey}

\subsection{Chronology}

Dates rest on the witnesses stated; where official acts and the
autobiography disagree (the Capua consecration), both are given.

\begin{lifetimeline}
\lifeevent{1542-10-04}{Birth at Montepulciano; parents Vincenzo
Bellarmino and Cinzia Cervini (sister of Marcellus II)}{CE 1907;
Lux illa (M, process-based)}{Secure}
\lifeevent{1560-09-20}{Enters the Society of Jesus at Rome; first
vows next day}{CE 1907; Lux illa}{Secure}
\lifeevent{1560--63}{Philosophy at the Roman College; master's
title}{Lux illa (M)}{Secure in outline}
\lifeevent{c.~1563--67}{Teaching at Florence, then Mondov\`i;
self-taught Greek; constant preaching}{autobiography (A); Lux
illa}{Secure in outline; year edges soft}
\lifeevent{1567--69}{Theology at Padua}{CE 1907; autobiography}
{Secure}
\lifeevent{1569}{Sent to Louvain via Milan, traveling with William
Allen}{CE 1907; autobiography (A)}{Secure}
\lifeevent{1570 (Lent--Easter)}{Orders in one season: tonsure,
minors, subdiaconate at Li\`ege; diaconate and priesthood at Ghent
from Bp.\ Cornelius Jansenius (Holy Saturday, 25 March); first
Mass at Louvain in the Easter octave}{autobiography (A);
Benedict XVI 2011 (M) for 25 March}{Secure; the one-season
compression is A-only}
\lifeevent{1570-10 -- 1576}{Teaches the Summa at Louvain, six
years; Hebrew; preaching}{autobiography (A); Providentissimus
Deus (M)}{Secure}
\lifeevent{1576}{Recalled to Italy; chair of Controversies at the
Roman College founded under Gregory XIII}{CE 1907; Lux illa}
{Secure; teaching-years count varies by witness (to 1586/87/88)}
\lifeevent{1586, 1588, 1593}{The three Ingolstadt tomes of the
\work{Disputationes}; third tome drafted at Frascati, autumn
1592}{Providentissimus Deus (M); autobiography (A)}{Secure}
\lifeevent{1588; 1592}{Spiritual father of the Roman College
(direction of Aloysius Gonzaga, d.\ 1591); rector}{CE 1907;
autobiography}{Secure}
\lifeevent{1590}{Legation to France with Cardinal Gaetani; siege
of Paris; Sixtus V places the \work{Controversiae} on his 1590
Index \latin{donec corrigantur}; Sixtus dies 27 August; listing
removed; return to Rome 11 November}{CE 1907; autobiography (A);
D\"ollinger--Reusch/Reusch Index bibliography (K)}{Secure; the
1590 Index never promulgated (K)}
\lifeevent{1591--92}{Vulgate revision counsel; Zagarolo
commission; Clementine Vulgate (1592) with Bellarmine's preface}
{autobiography (A); CE 1907}{Secure; the preface's morality
debated then and since}
\lifeevent{1594--c.\,96}{Provincial of Naples}{Lux illa (M)}
{Secure}
\lifeevent{c.~1597}{Recalled by Clement VIII; papal theologian;
consultor of the Holy Office; \work{Dottrina cristiana breve}
(1597) and \work{Dichiarazione} (1598)}{Lux illa; CE 1907;
Benedict XVI 2011}{Secure}
\lifeevent{1599-03-03}{Cardinal priest of S.~Maria in Via;
``non habet parem Ecclesia Dei quod ad doctrinam''}{Lux illa (M);
CE 1907}{Secure}
\lifeevent{1598--1607}{\latin{De auxiliis} controversy: commission
1598; disputations before Clement VIII 1602--05 and Paul V
1605--07; Bellarmine counsels against a decision; Paul V ends it
5 September 1607 undecided}{CE 1908 ``Congregatio de Auxiliis''
(K/E); CE 1907}{Sequence secure at reported level}
\lifeevent{1600-02-17}{Bruno burned; Bellarmine one of the
cardinal inquisitors; his individual role undocumented in
witnesses read here}{CE 1908 ``Giordano Bruno'' (K); Opere XIX
322 for his Inquisition membership (N)}{Membership secure;
personal role open}
\lifeevent{1602}{Named to Capua (18 March, per Benedict XVI);
consecrated by Clement VIII on Good Shepherd Sunday = 21 April
(A; Benedict XVI) or 20 April (Lux illa); enters Capua 1 May}
{autobiography (A); Lux illa (M); Benedict XVI (M)}{Appointment
secure; consecration day: 21 April preferred, conflict recorded}
\lifeevent{1602--05}{Residence at Capua: three visitations, three
diocesan synods, one provincial council; organized monthly relief
of the poor}{autobiography (A); Roma altera (M)}{Secure}
\lifeevent{1605}{Conclaves of Leo XI and Paul V; ``little was
lacking that he should be made Pope'' (his own words); resigns
Capua; retained at Rome}{autobiography (A); CE 1907}{Discussion
of his candidacy secure; nearness not measurable}
\lifeevent{1606--07}{Venetian interdict pamphlet war: Bellarmine
and Baronius against Sarpi and Marsilio}{CE 1907 (E/K);
autobiography (A, ``two or three little Italian books'')}{Secure
in outline}
\lifeevent{1606--10}{Oath of Allegiance exchange: breves of 1606
and 1607; letter to Blackwell (28 September 1607);
\work{Triplici nodo} (1607); \work{Responsio Matthaei Torti}
(1608); \work{Premonition} and \work{Apologia} (1609); against
Barclay, \work{De potestate Summi Pontificis} (1610), condemned
by the Parlement of Paris}{McIlwain texts (N); CE 1907}{Secure}
\lifeevent{1611-04}{Galileo's telescope findings queried with the
Roman College mathematicians (19 April); their confirmation
(24 April)}{Opere XI, nos.\ 515, 520 (A/N)}{Secure}
\lifeevent{1615-04-12}{Letter to Foscarini: \latin{ex
suppositione} tolerated; assertion dangerous; demonstration
would compel reinterpretation; no demonstration believed}{Opere
XII, no.\ 1110, contemporary copy (A)}{Text secure at copy
level}
\lifeevent{1616-02/03}{Qualifiers' censure (24 Feb, without
Bellarmine); papal direction (25 Feb); warning and unsigned
precept minute (26 Feb); Bellarmine's report of acquiescence
(3 Mar); Index decree (5 Mar), Copernicus suspended \latin{donec
corrigatur}, Foscarini condemned}{Opere XIX, 320--323, 278 (N)}
{Documents secure; course of 26 Feb unresolved}
\lifeevent{1616-05-26}{Certificate to Galileo in Bellarmine's
hand: no abjuration, no penance; the doctrine ``may not be
defended or held''}{Opere XIX, 348 (autograph) and 342 (Galileo's
copy) (A)}{Secure}
\lifeevent{1611--20}{The late works: Psalms commentary (1611);
\work{De scriptoribus ecclesiasticis} (1612); autobiography
(1613); the five retreat books (1615--1620)}{CE 1907;
Providentissimus Deus (M); autobiography (A)}{Secure}
\lifeevent{1621-09-17}{Death at Sant'Andrea al Quirinale, aged
78; popular acclamation; burial in the Ges\`u}{Lux illa (M,
process-based); CE 1907}{Secure}
\lifeevent{1622--27}{Ordinary processes (Rome, Montepulciano,
Capua, Naples); Urban VIII signs the commission 12 September
1626; introduction of the cause 15 January 1627}{Lux illa (M);
D\"ollinger--Reusch (K)}{Secure}
\lifeevent{1675--77}{Writings cleared (1674); virtues voted
favorably by consultors (7 Sept 1675); blocked in the general
congregation of 26 Sept 1677 under Innocent XI (6 cardinals
against, incl.\ Barbarigo, Azzolini, Casanate)}
{D\"ollinger--Reusch, from the printed vota (K)}{Secure at
reported level}
\lifeevent{1711--19; 1748--53}{Clement XI resumption (Lambertini
promoter); writings recleared 1714--19; Benedict XIV resumption
1748; favorable vote 5 May 1753 unconfirmed; Passionei's votum
and the French court's intervention; Benedict XIV's ``sad
circumstances of the times''}{D\"ollinger--Reusch (K); CE 1907
(E/K)}{Secure at reported level}
\lifeevent{1920--23}{Heroic virtues decreed 22 December 1920;
miracles approved 15 April 1923; beatified 13 May 1923; remains
translated to Sant'Ignazio 21 June 1923}{Lux illa (M)}{Secure}
\lifeevent{1930-06-29}{Canonized with the North American Martyrs
and Theophilus of Corte; memorial assigned to 17 September}{Lux
illa, AAS 22, 593--604 (M)}{Secure}
\lifeevent{1931-09-17}{Declared Doctor of the Universal Church,
\work{Providentissimus Deus}; feast then on 13 May extended to
the whole Church}{AAS 23, 433--438 (M)}{Secure}
\lifeevent{1967-09-29}{Co-patron, with St.~Stephen Protomartyr,
of the archdiocese of Capua (\work{Roma altera})}{AAS 60, 12--13
(M)}{Secure; object limited to Capua}
\lifeevent{current}{Optional memorial, 17 September, General
Roman Calendar}{2002 \work{Missale Romanum} calendar (M)}{Secure
as of 2026-07-25}
\end{lifetimeline}

\subsection{Method, limits, and negative results}

Earliest controlled evidence governs; the autobiography outranks
the official retrospectives for lived detail but is discounted as
self-presentation; institutional records are read for what they
are. Linguistic limits: Latin and Italian sources were read
directly in public scans; German (D\"ollinger--Reusch's
commentary) likewise; no archival document, manuscript, or
critical edition beyond those scans was examined. Bounded negative
results from the witnesses read (all retrieved 2026-07-25): no
witness read here names Bellarmine's individual acts or votes in
the Bruno case; no early witness was located for the
``his opponents' arguments stated better'' quip; no competent
dated act was located for patronage of catechists, catechumens,
or canonists; no incorruption claim was located; the 26~February
1616 minute bears no signatures, and no witness resolves its
course of events; no conclave voting record was inspected for the
1605 near-election. These are bounded, correctable results of the
searches described in the research records, not proofs of
absence. Mutable online sources (vatican.va, newadvent.org,
archive.org) are cited as retrieved 2026-07-25.

\subsection{Rights and review state}

Project prose is original. Quotations follow public-domain
sources: the Latin and Italian originals in D\"ollinger--Reusch
(1887), Favaro's edition (1890--1909), and the early printings;
the Sturge translation of von Gebler (1879); McIlwain's Harvard
edition (1918) for the English texts of James's side; the
1907--1912 Catholic Encyclopedia. Brief quotations from the
\work{Acta Apostolicae Sedis} and from Benedict XVI's catechesis
are cited with attribution for reporting and study; project
English renderings of Latin and Italian sources are marked as
project translations. This publication has received internal
editorial and production checking only: no independent
historical, theological, canonical, archival, scientific, or
rights review, and no ecclesiastical approval, is claimed or
implied.
